% ------------------------------- Chapter 4 -------------------------------
\chapter{Rotor Blade Construction and CFD Evaluation}

The blade profile obtained in Chapter~3 is assembled into a linear cascade and
evaluated with a 2-D RANS computation. This chapter describes the passage
geometry, the mesh, the solver setup and convergence, and the post-processed
aerodynamic results (loading, Mach/pressure fields and the exit wake).

% =====================================================================
\section{Blade Passage Geometry}

The cascade is modelled as three blades: the central (measured) blade and its two
neighbours, whose facing surfaces form the upper and lower passage walls. The
non-dimensional pitch and throat characterise the passage:
\begin{equation}
    s/c_x = 1.136,\quad o/c_x = 0.805,\quad o/s = 0.709,\quad
    (x/c_x)_{throat} = 0.284,
\end{equation}
with dimensional values $c_x = \SI{0.149}{m}$, $s = \SI{0.170}{m}$ and throat
$o = \SI{0.120}{m}$, for $N_b = 42$ blades.

\section{Evaluation of Profile Curvature and Passage Area}

The prescribed thickness distribution and the analytic curvature of the surfaces
are shown in Fig.~\ref{fig:thkcur}; the passage-width (area) variation along the
axial direction, which contracts monotonically to the throat, is shown in
Fig.~\ref{fig:passwidth}. A smooth, monotonic contraction without local
diffusion is desirable to avoid boundary-layer growth on the suction side.


\begin{figure}[H]
\centering
\includegraphics[width=0.7\textwidth]{Passage Width Variation.png}
\caption{Passage-width (area) variation through the blade row.}
\label{fig:passwidth}
\end{figure}

\begin{figure}[H]
\centering
\begin{subfigure}{0.3\textwidth}
    \includegraphics[width=\textwidth]{Blade Cascade.png}
    \caption{Cascade of rotor blades.}
\end{subfigure}\hfill
\begin{subfigure}{0.4\textwidth}
    \includegraphics[width=\textwidth]{Rotor Blade Passage Geometry.png}
    \caption{Blade-to-blade passage and throat.}
\end{subfigure}
\caption{Rotor-blade cascade and passage geometry.}
\label{fig:passage}
\end{figure}

% =====================================================================

\begin{figure}[H]
\centering
\begin{subfigure}{0.48\textwidth}
    \includegraphics[width=\textwidth]{Thickness Distribution.png}
    \caption{Thickness distribution.}
\end{subfigure}\hfill
\begin{subfigure}{0.48\textwidth}
    \includegraphics[width=\textwidth]{Curvature Distribution.png}
    \caption{Curvature $\kappa$ distribution.}
\end{subfigure}
\caption{Profile thickness and curvature along the normalised surface coordinate.}
\label{fig:thkcur}
\end{figure}



% =====================================================================
\section{CFD Evaluation}

\subsection{Mesh Generation}

An unstructured triangular mesh is generated with GMSH~\cite{gmsh} over the
three-blade passage. The domain is bounded by the central blade (no-slip wall),
the two neighbour walls (the pressure side of the upper blade and the suction side
of the lower blade), and inlet/outlet planes extended half an axial chord upstream
and downstream along the flow angles $+\beta_2$ and $-\beta_3$. The inlet and
outlet heights are both $\approx 2s = \SI{339.2}{mm}$. A distance-based size field
refines the cells near all blade walls. The mesh contains about $1.81\times10^4$
nodes and $3.54\times10^4$ elements.


\begin{figure}[H]
\centering
\begin{subfigure}{0.35\textwidth}
    \includegraphics[width=\textwidth]{Mesh-geo-verification.png}
    \caption{Geometry verification of the 3-blade domain.}
\end{subfigure}\hfill
\begin{subfigure}{0.65\textwidth}
    \includegraphics[width=\textwidth]{Blade-mesh.png}
    \caption{Generated GMSH mesh.}
\end{subfigure}
\caption{Three-blade cascade domain and mesh.}
\label{fig:mesh}
\end{figure}


\subsection{CFD Setting}

The flow is solved with SU2~\cite{su2} (compressible RANS, Spalart--Allmaras
turbulence model). The boundary conditions are imposed in the rotor-relative
frame: stagnation inlet from station~2 ($T_{02R}$, $P_{02R}$, flow direction
$\beta_2$), static back-pressure $P_3$ at the outlet, a no-slip adiabatic central
blade and inviscid (Euler) neighbour walls. The main settings are listed in
Table~\ref{tab:su2}.


\begin{table}[H]
\centering
\caption{SU2 CFD solver configuration.}
\label{tab:su2}
\begin{tabular}{l l}
\toprule
\textbf{Setting} & \textbf{Value}\\
\midrule
Solver                       & RANS (compressible) \\
Turbulence model             & Spalart--Allmaras \\
Inlet total temperature $T_{02R}$ [K] & 1578.1 \\
Inlet total pressure $P_{02R}$ [Pa]   & 1521338 \\
Inlet flow angle $\beta_2$ [\degr]    & 42.64 \\
Outlet static pressure $P_3$ [Pa]     & 1107265 \\
Inlet relative Mach $M_{w2}$          & 0.308 \\
Central blade (WALL\_MID)    & no-slip, adiabatic \\
Neighbour walls (WALL\_UP/LO)& Euler (slip) \\
Convergence monitor          & RMS density residual \\
\bottomrule
\end{tabular}
\end{table}

\subsection{CFD Convergence}

The simulation is advanced until the density residual drops by several orders of
magnitude (RMS$[\rho]$ falling below $-7.7$), at which point the surface loading
and the integral exit quantities are stationary. The solution is then used for
the post-processing below.

% =====================================================================
\section{Isentropic Static-Pressure Distribution}

The blade loading is examined through the isentropic Mach number and the static
pressure on the suction and pressure surfaces (Fig.~\ref{fig:loading}). The
suction-side isentropic Mach number rises to $M_{is}\approx 0.84$ near the throat
before recovering, while the pressure side stays below $\approx 0.79$; the static
pressure spans $1.02$--$\SI{1.53}{MPa}$ on the pressure side and
$0.98$--$\SI{1.52}{MPa}$ on the suction side. The loading is front-to-mid loaded
with no sign of suction-side separation.

\begin{figure}[H]
\centering
\begin{subfigure}{0.48\textwidth}
    \includegraphics[width=\textwidth]{Isentropic_Mach.png}
    \caption{Isentropic Mach number.}
\end{subfigure}\hfill
\begin{subfigure}{0.48\textwidth}
    \includegraphics[width=\textwidth]{Pressure Distribution.png}
    \caption{Static pressure on PS / SS.}
\end{subfigure}
\caption{Blade loading: isentropic Mach and static-pressure distribution.}
\label{fig:loading}
\end{figure}

% =====================================================================
\section{Colourmaps of the CFD Analysis}

The Mach-number and static-pressure fields (Fig.~\ref{fig:colourmaps}) show the
flow accelerating through the converging passage, the peak velocity on the
suction side near the throat, and the smooth pressure relaxation towards the
outlet, consistent with the prescribed expansion to $P_3$.

\begin{figure}[H]
\centering
\includegraphics[width=0.95\textwidth]{CFD_Colormaps.png}
\caption{CFD colourmaps of the cascade flow: Mach number, specific total
enthalpy, entropy rise and static pressure.}
\label{fig:colourmaps}
\end{figure}

% =====================================================================
\section{Velocity Profile and Wake}

The pitchwise distributions one trailing-edge region downstream
(Fig.~\ref{fig:wake}) reveal the blade wake as a momentum deficit in all four
quantities -- total pressure, Mach number, flow angle and axial velocity. At the
extraction plane $x = x_{TE}+0.30\,c_x = \SI{0.194}{m}$ the mass-averaged total
pressure is $P_t = \SI{1.49}{MPa}$ ($P_t/P_{t0}=0.978$), the mean Mach number is
$0.681$ and the mean flow angle is $-70.7\degr$. The axial velocity drops from
$\approx \SI{191}{m/s}$ in the free stream to $\approx \SI{158}{m/s}$ in the wake
core (a $9.2\%$ deficit), narrow and well-mixed.

\begin{figure}[H]
\centering
\includegraphics[width=0.88\textwidth]{Pitchwise_Distributions.png}
\caption{Exit-plane pitchwise distributions of total pressure, Mach number, flow
angle and axial velocity (wake), shown un-normalised.}
\label{fig:wake}
\end{figure}

% =====================================================================
\section{CFD Prediction Result}

The CFD confirms the mean-line design: the relative exit Mach number
($M_{w3}\approx 0.67$--$0.68$) and exit angle ($\approx 70\degr$ relative) match
the design intent, the passage delivers the prescribed expansion to $P_3$, and
the stage total-pressure loss across the rotor is modest ($P_t/P_{t0}=0.978$,
i.e.\ $\approx 2.2\%$) with a narrow, attached wake. The blade row is therefore
aerodynamically sound for the chosen mean-line operating point.
